% Copyright (c) 2026 Geovana Neves. Licensed under CC BY 4.0 (https://creativecommons.org/licenses/by/4.0/). See LICENSE-AND-AUTHORSHIP.md.
%
% 00-fts-overview/text/05-guides.tex
%
% Part 5: the map of the other seven guides, the order to read them in for
% four kinds of reader, the public geometries to practise on, and how to cite
% the tool.

\section{The guides}

\begin{frame}{Eight guides, and when to read each}
\relax
{\ftstablesize\renewcommand{\arraystretch}{1.08}
\begin{tabularx}{\linewidth}{@{}P{0.04\linewidth}P{0.31\linewidth}RP{0.24\linewidth}@{}}
\toprule
& \textbf{File} \code{pyfts-guide-} & \textbf{What it covers} & \textbf{Read it when} \\
\midrule
0 & \code{00-fts-overview} & the solver, the workflow, the map &
  before anything else \\
1 & \code{01-workspaces} & the tree, the input files, choosing a
  workflow, running, what a run leaves, two workspaces and sync, version
  control & you set up or run a study \\
2 & \code{02-gui-to-pyfs} & every step of a GUI session, and the
  key that takes it & you know the interface \\
3 & \code{03-references} & lengths, the moment point, aliases,
  frames, rotors, actuator discs & you write a reference file \\
4 & \code{04-solver-setup} & the setup file, its keys and the
  solver commands behind them & you tune how a point is solved \\
5 & \code{05-pproc-definitions} & what every product, reduction
  and window means & you read or plot a result \\
6 & \code{06-fsi} & the fluid-structure loop, the structural
  input, calibration, limits [22] & a structure deforms under load \\
7 & \code{07-python-environment-offline} & Python, a virtual
  environment and the package with no internet & once per offline machine \\
\bottomrule
\end{tabularx}}
\end{frame}

\begin{frame}{Four ways through them}
\begin{columns}[T]
\begin{column}{0.48\textwidth}
\begin{block}{A first study}
\small 0, then 1 (Parts 1 to 11), then 3 and 4 for the files you cite, then 5
before you plot. The package page ``Getting started'' is the shortest path to
a first run [21].
\end{block}
\begin{block}{From the interface to the workspace}
\small 0 (Parts 3 and 4), then 2 side by side with your own session, then 4.
\end{block}
\end{column}
\begin{column}{0.48\textwidth}
\begin{block}{A rotor}
\small 0 (Part 2, rotation), 3 (the rotor block), 1 (choosing a workflow and
the quasi-steady rotor), 5 (per-blade, phase-locked and window products); 6
for a blade that deforms.
\end{block}
\begin{block}{Installing and operating}
\small 7 on each machine without internet, then 1 (the cluster, two
workspaces and sync, space on disk).
\end{block}
\end{column}
\end{columns}
\begin{takeaway}Each guide is written for one version of the package; its
title page says which.\end{takeaway}
\end{frame}

\begin{frame}{Where to practise, and how to cite}
\begin{columns}[T]
\begin{column}{0.48\textwidth}
\begin{block}{Two public geometries}
\begin{itemize}\small
  \item \textbf{SMI}: an airframe geometry, wing and wing-body, published
        for practice [18].
  \item \textbf{TUD-XPROP}: a research propeller of Delft University of
        Technology, requested from its own record and cited by it [19].
\end{itemize}
\end{block}
\begin{block}{Where to take a question}
\small Solver behaviour to the vendor, with the build and the page. The
driver to the repository, as an issue.
\end{block}
\end{column}
\begin{column}{0.48\textwidth}
\begin{block}{Citing the tool}
\small \pkgcitation
\end{block}
\begin{block}{The Python library}
\small For scripts written in Python rather than a matrix: the library
guide, \code{pyflightstream\_user\_guide.tex} in the repository's
\code{guide/} folder.
\end{block}
\end{column}
\end{columns}
\end{frame}
